\subsection{Delegating Conjunctive Queries}
Under this query execution model, we now show that correct delegation to composite
sources can be established under \emph{incomplete} discovery.

\begin{lemma}[Claimed triples occur only in their claimant]
\label{lem:exclusive}
Under Assumption~\ref{as:disjoint}, let $c \in \Dd$ claim $\tau$. Then $\tau$ is
admitted from no source $d \neq c$. Consequently, for every $A \subseteq \Dd$,
$\tau \in G^{\alpha}_A$ implies $c \in A$ and $\tau \in \dataf(c)$.
\end{lemma}

\begin{proof}
Suppose $\tau$ is admitted from some $d \neq c$. The second clause of
\eqref{eq:admitted} fails, as $c \in \Dd \setminus \{d\}$ claims $\tau$. So the first
holds and $d$ claims $\tau$, contradicting Assumption~\ref{as:disjoint}.
\end{proof}

A delegated composite answers over all of $\dataf(c)$, so its answer must be filtered
to what the claim semantics admits. The engine can always check the claim of $c$, as
$\alpha_c$ is known once $c$ is retrieved. For a BGP $B'$ we therefore keep the solutions
whose matched triples $c$ all claims,
\[
  \Phi_c(\Omega) \;=\; \bigl\{\, \mu \in \Omega \;\bigm|\; c \text{ claims } \mu(t)
  \text{ for every } t \in B' \,\bigr\},
\]
and evaluate the rest of $B'$ over the other sources. This refines
Definition~\ref{def:delegation} for $S = \{c\}$:
\[
  R \;:=\; \Phi_c\bigl(\ev{B'}{\dataf(c)}\bigr) \;\cup\;
  \ev{B'}{G^{\alpha}_{\Dd \setminus \{c\}}}.
\]
Unclaimed triples held by $c$ are also held by the sources it aggregates, so they are
found by the second term, provided these sources are reachable.

\begin{assumption}[Reachable aggregation]
\label{as:agg-reach}
If a composite $c$ is in $\Dd$, so is every source it aggregates: $\Dd_c \subseteq \Dd$.
\end{assumption}

This holds, for instance, when a composite exposes its sources and $\disc$ follows them.
What can go wrong is a solution that the two terms split between them.

\begin{definition}[Claim-uniform sub-query]
\label{def:uniform}
Let $G = G^{\alpha}_{\Dd}$. A BGP $B'$ is \emph{claim-uniform} for a composite $c$ if
every $\mu \in \ev{B'}{G}$ has $c$ claim either all of $\mu(t)$, $t \in B'$, or none of
them.
\end{definition}

\begin{theorem}[Delegation to a claim]
\label{th:claim-delegation}
Under Assumptions~\ref{as:disjoint} and~\ref{as:agg-reach}, let $B' \subseteq B$ be
non-empty, $c \in \Dd(T)$ a composite, and $G = G^{\alpha}_{\Dd}$. Then
$R = \ev{B'}{G}$ if and only if $B'$ is claim-uniform for $c$; in that case delegating
$B'$ to $c$ is correct. Consequently, whenever claim uniformity follows from $B'$ and
$\alpha_c$ alone, correctness is determined at $T$, whether or not $F(T) = \emptyset$.
\end{theorem}

\begin{proof}
Write $G_c := G^{\alpha}_{\{c\}}$ and $G_{\neg c} := G^{\alpha}_{\Dd \setminus \{c\}}$.

($R \subseteq \ev{B'}{G}$.) A triple of $\dataf(c)$ that $c$ claims is admitted from
$c$ by the first clause of \eqref{eq:admitted}, so a solution kept by $\Phi_c$ maps into
$G_c \subseteq G$. The second term maps into $G_{\neg c} \subseteq G$.

($\ev{B'}{G} \subseteq R$ under uniformity.) Let $\mu \in \ev{B'}{G}$. If $c$ claims every
$\mu(t)$, each is admitted only from $c$ by Lemma~\ref{lem:exclusive}, hence lies in
$\dataf(c)$; so $\mu \in \ev{B'}{\dataf(c)}$ and $\Phi_c$ keeps it. If $c$ claims no
$\mu(t)$, let $\tau = \mu(t)$ be admitted from $d$. If $d \neq c$, then
$\tau \in G_{\neg c}$. If $d = c$, then $\tau$ is unclaimed, as $c$ does not claim it and
by \eqref{eq:admitted} no other source does; it is held by some $d' \in \Dd_c$, which lies
in $\Dd$ by Assumption~\ref{as:agg-reach} and admits $\tau$, so again
$\tau \in G_{\neg c}$. Hence $\mu \in \ev{B'}{G_{\neg c}}$.

(Necessity.) Let $\mu \in \ev{B'}{G}$ mix claimed and unclaimed triples, with $c$
claiming $\mu(t_1)$ but not $\mu(t_2)$. Then $\Phi_c$ discards $\mu$, and
$\mu(t_1) \notin G_{\neg c}$ by Lemma~\ref{lem:exclusive}, so $\mu \notin R$.

Given $R = \ev{B'}{G}$, $R \bowtie \ev{B \setminus B'}{G} = \ev{B}{G}$, which is
correctness. Finally, $\Phi_c$ depends only on $\alpha_c$, and the argument quantifies over
every source of $\Dd$, discovered or not; so whenever uniformity is guaranteed by $B'$ and
$\alpha_c$ alone, the verdict is the same in every state observationally equivalent to
$W(T)$.
\end{proof}

Theorem~\ref{th:claim-delegation} pins down what remains of the adversary of
Theorem~\ref{th:undetermined}: it wins only by building a solution that mixes a triple
claimed by $c$ with one that is not, so that the join crosses the boundary of $c$. Claim
uniformity is exactly the absence of such solutions. A triple placed behind the frontier
that $c$ claims is admitted only from $c$, and one that $c$ does not claim is found by the
second term; neither breaks a uniform sub-query.

The simplest uniform case needs no knowledge of the data at all: if $c$ claims every
triple matching every pattern of $B'$, all solutions are fully claimed. We say $B'$ is
then \emph{within the claim} of $c$; the second term of $R$ is empty and
$\Phi_c$ is the identity. Under namespace claims (Section~\ref{sec:namespace-claims}) this
requires constant subjects in the namespace of $c$. Uniformity is much weaker: it only
forbids mixed solutions, and Section~\ref{sec:shapes} shows that subject-determined
claims guarantee it for star-shaped sub-queries regardless of the data. A composite that
claims nothing has $\Phi_c(\Omega) = \emptyset$, so delegating to it is vacuous: we
delegate only to claimants, and a single document only as a degenerate composite that
declares a claim.


\todo{Lemma on how theorem 1 holds if we filter any sources in scope 
of $c$ both backwards and forwards in time. Mention how the filtering changes
nothing given the assumptions / semantics of our model}
\begin{proof}
  Sketch
  Part 1. If we filter -> Theorem~\ref{th:correct-delegation} can be directly applied
  Part 2. Filtering is possible: forward -> easy just take $\disc()$ 
  and apply setminus of in scope documents.
  Backwards -> deduplication, we only have duplicate results. So filter any
  intermediate results at time T in data(scope) and filter any produced results
  in composite resource on B prime
\end{proof}

